\begin{ExerciseBlock}
\ExerciseHeading{05}{Massen am Tisch}{work_and_energy_masses_and_pulleys}
Auf einem Holztisch wird eine Masse $m$ aufgrund der Gewichtskräfte der Massen $3m$ und $6m$ bewegt (siehe Abbildung). Dabei tritt eine Gleitreibung mit Koeffizienten $\mu_G$ auf. Anfangs befinden sich die Massen in Ruhe, wobei die auf dem Tisch aufliegende Masse $m$ einen Abstand $s$ zur rechten Rolle hat.

\begin{center}
	\CachedTikzFigure{content/resources/chapters/work_and_energy/tikz/masses_and_pulleys.tex}
\end{center}

Zeigen Sie, dass die Masse $m$ mit der Geschwindigkeit \[v = \sqrt{\frac{gs}{5} \left( 3 - \mu_G\right)} \] auf die Rolle trifft. \\
Vernachlässigen Sie in der Rechnung die Rotation der Rollen.
\end{ExerciseBlock}
\UseExerciseSolution{work_and_energy_masses_and_pulleys}
